%==============================================================================
% SUPPORTING INFORMATION FOR JOURNAL OF CHEMICAL THEORY AND COMPUTATION
%==============================================================================
\documentclass[11pt,a4paper]{article}
\usepackage[utf8]{inputenc}
\usepackage[T1]{fontenc}
\usepackage{amsmath,amssymb,amsfonts,graphicx,booktabs,siunitx,physics,bm}
\usepackage{xcolor}
\usepackage{hyperref}
\usepackage[margin=2.5cm]{geometry}

% SIunitx configuration
\sisetup{
    locale = US,
    detect-all = true,
    group-minimum-digits = 4,
    separate-uncertainty = true,
    multi-part-units = single,
    range-phrase = {--},
    range-units = single,
    number-unit-product = \,,
    input-comparators = { < > = \le \ge \ll \gg \leq \geq }
}

\begin{document}

\title{\textbf{Supporting Information for:\\Scalable Simulation of Non-Markovian Energy Transfer via Process Tensor Hierarchy of Pure States with Stochastically Bundled Dissipators}}

\author{Steve Cabrel Teguia Kouam, Theodore Goumai Vedekoi, Jean-Pierre Tchapet Njafa,\\ Jean-Pierre Nguenang, and Serge Guy Nana Engo}

\date{\today}
\maketitle

%==============================================================================
% SECTION S1: FMO HAMILTONIAN PARAMETERS
%==============================================================================
\section{FMO Hamiltonian and Spectral Density Parameters}

While the PT-HOPS/SBD formalism and 12-test validation benchmarking are documented comprehensively in the main text, this Supporting Information explicitly provides the 7-site bacteriochlorophyll-\textit{a} matrices utilized in our numerical models (from Ref.\ \cite{Adolphs2006}).

\subsection{Site Energies}

\begin{table}[h]
\centering
\caption{\textbf{FMO Site Energies (cm$^{-1}$)}}
\label{tab:SI_site_energies}
\begin{tabular}{ccccccc}
\toprule
$\varepsilon_1$ & $\varepsilon_2$ & $\varepsilon_3$ & $\varepsilon_4$ & $\varepsilon_5$ & $\varepsilon_6$ & $\varepsilon_7$ \\
\midrule
12400 & 12550 & 12250 & 12350 & 12450 & 12600 & 12500 \\
\bottomrule
\end{tabular}
\end{table}

\subsection{Electronic Couplings}

The inter-chromophore electronic coupling matrix $J_{nm}$ (in cm$^{-1}$):

\begin{equation}
\mathbf{J} = \begin{pmatrix}
0 & -87.7 & 5.5 & -5.9 & 6.7 & -13.7 & -9.9 \\
-87.7 & 0 & 30.8 & 8.2 & 0.7 & -11.9 & 4.3 \\
5.5 & 30.8 & 0 & -53.5 & -2.2 & -9.6 & 6.0 \\
-5.9 & 8.2 & -53.5 & 0 & 1.2 & -6.3 & 1.7 \\
6.7 & 0.7 & -2.2 & 1.2 & 0 & 9.6 & -6.0 \\
-13.7 & -11.9 & -9.6 & -6.3 & 9.6 & 0 & 89.7 \\
-9.9 & 4.3 & 6.0 & 1.7 & -6.0 & 89.7 & 0
\end{pmatrix}.
\label{eq:SI_coupling_matrix}
\end{equation}

\subsection{Spectral Density Parameters}

\begin{table}[h]
\centering
\caption{\textbf{Spectral Density Parameters}}
\label{tab:SI_spectral_params}
\begin{tabular}{lc}
\toprule
\textbf{Parameter} & \textbf{Value} \\
\midrule
Drude--Lorentz reorganization ($\lambda_{\rm DL}$) & \SI{35}{\per\centi\meter} \\
Drude--Lorentz cutoff ($\gamma_{\rm DL}$) & \SI{50}{\per\centi\meter} \\
Vibronic mode 1 ($\omega_1$) & \SI{150}{\per\centi\meter} \\
Vibronic mode 2 ($\omega_2$) & \SI{200}{\per\centi\meter} \\
Vibronic mode 3 ($\omega_3$) & \SI{575}{\per\centi\meter} \\
Vibronic mode 4 ($\omega_4$) & \SI{1185}{\per\centi\meter} \\
Huang--Rhys $S_1$ & 0.05 \\
Huang--Rhys $S_2$ & 0.02 \\
Huang--Rhys $S_3$ & 0.01 \\
Huang--Rhys $S_4$ & 0.005 \\
\bottomrule
\end{tabular}
\end{table}

%==============================================================================
% REFERENCES
%==============================================================================
\begin{thebibliography}{99}
\bibitem{Adolphs2006}
Adolphs, J.; Renger, T. Biophys. J. \textbf{2006}, \textit{91}, 2778--2797.
\end{thebibliography}

\end{document}
